
\section{Experiments}
\label{sec:experiments}

In this section, we describe in detail the experimental setup and the experiments conducted to compare our methods to baselines. Our experiments include a large number of simulated experiments to validate the algorithms, and a small number of real world experiments. \new{In the simulated experiments, we use a 6-dof UR5 mounted on a 3-\dof planar mobile platform, yielding in a 9-\dof robotic system. In the physical experiments, we use a fixed-base 6-\dof UR5, see Fig. \ref{fig:physical:experiments:two:2}a. Although the physical arm is mounted on a mobile base, the platform was kept stationary throughout the real-world experiments.} \remove{In all simulated and real-world experiments, we use a UR5 robotic arm mounted on a Clearpath Husky A200 mobile platform. However, the base remained stationary throughout all experimental trials.} A flashlight attached to the arm's end-effector served as our visibility-based instrument. We assign a range of $h=3$~m to the flashlight, and aperture angles of $\gamma\approx30^\circ$ and $\gamma\approx10^\circ$ in the simulated and real-world experiments, respectively.%  

To ensure consistency across evaluations, targets were pre\-/instantiated and shared across all algorithms. Internal algorithmic stochasticity, e.g., sampling subroutines, remained independent. Since randomly generated instances $(r, E, s, P_S)$ do not guarantee the existence of an admissible solution, all reported runtimes correspond exclusively to successful trials unless otherwise stated. In all of the experiments planning was done on a desktop PC with an Intel i7-14700F CPU, 64 GB of RAM, and a GeForce RTX 4090 GPU. \new{A video illustrating the simulated and real-world experiments is included in the supplementary materials.}

% ---------------------------------------------------

\subsection{Algorithms evaluation}
\label{sec:experiments:sim}

\new{The proposed VisTAMP algorithms were evaluated in the simulation shown in Figure \ref{fig:simulated:experiments}. The environment comprises two central rooms encircled by a wide corridor. The first room is highly constrained, featuring a narrow passage flanked by ten deep, open-top containers (five per side) that restrict base rotation and severely limit the configurations capable of achieving line-of-sight to the container floors. The second room is an unconstrained workspace with three windows per outer wall and three shelves per inner wall. Being obstacle-free, its valid configurations form a large connected component in \cspace. Spherical targets, whose radii match their clearance to the nearest obstacle, were randomly sampled within the containers or on the shelves.}

% \new{ We first evaluate the performance of the proposed \vistamp algorithms in a simulated environment, shown in Figure \ref{fig:simulated:experiments}. Two rooms sharing a wall are located at the center of the environment, surrounded by a wide path. One room is windowless, and has a narrow doorway in one of its walls, leading to a same-width corridor passing between 10 deep open-top containers, 5 on each side. The width of the passage allows the robot to enter and move in a very limited range of base angles, prohibiting a full rotation inside. and the depth of the containers means that only a small subset of manipulator configurations can place targets at their bottom within the \fov. 
% \\
% The second room is door-less, and has three windows on each of its outer walls, and three shelves are mounted on each of its interior walls. The interior of this room is clear of obstacles, and configurations of the robot inside the room comprise a large connected component of the robot's \cspace.
% \\
% Spherical visibility targets were generated by randomly choosing a point in one of the containers or on one of the shelves, and choosing the clearance of that point, i.e., distance to the closest obstacle, as its radius.}


\new{We evaluate the performance of \visprm and \visrrt by comparing them against \tvmp adaptations of \prm, \rrt, and an adaptation of the \vir algorithm.} The baseline \prm modification is drawn from \cite{mrph-ocpuvsd-21}, which addresses a surface-covering problem for UV disinfection. In this framework, a roadmap $G=(V,E)$ is constructed using uniform sampling, and given a query $(s,P_S)$, $s$ is connected to $G$, and $V$ is scanned for a node $v$ such that $P_S$ is sufficiently covered by $\visatcfg{v}$. If no such configuration is found, the roadmap is expanded, and the \fov of configurations is evaluated as possible goals as they are added. In the case of the baseline \rrt, we simply check the \fov of every new configuration added, to determine if it satisfies the visibility requirements for $P_S$. Because these baselines represent minimal adaptations for the \tvmp problem, we retain their names to avoid introducing redundant nomenclature. \new{The original \vir algorithm \cite{zhvml-c3dfrvuvi-19} assumes a simple point robot and cannot be applied directly to complex manipulators. To adapt it as a benchmark, we have modified the query structure, evaluate cluster visibility relationships, and use retrieved points to compute robot configurations. We use its random sampling variant, as the grid approach yields preprocessing runtimes orders of magnitude worse than our proposed method.}

% \new{The original \vir algorithm \cite{zhvml-c3dfrvuvi-19} is a \prm-based method that computes the pairwise visibility of a set of of 3D workspace points, either a regular grid or a randomly sampled set, and then clusters these into sets of high \vi score. Given a visibility target, the cluster containing the target is identified, and samples from the cluster are taken as candidate configurations to be added to the roadmap. We adjust the  \vir for our problem by computing visibility relationships between clusters and using a modified query. These changes are required to utilize the algorithm with robots more complex than the point robot assumed in \cite{zhvml-c3dfrvuvi-19}. Our implementation uses a randomly sampled set of points, as the grid variant produced preprocessing runtimes several orders of magnitude worse than the \vi-tree even for coarse grid resolutions.}

%%%%%%%%%%%%%%%%%%%%%%%%%%%%%
% \begin{figure}[t!]
%     \centering
%     \begin{tabular}{ccc}
%         \includegraphics[trim={35mm 0mm 35mm 40mm}, clip, height=3.8cm]{walls_1.png} &~~~~~&  \includegraphics[trim={35mm 5mm 35mm 35mm}, clip, height=3.8cm]{walls_3.png}\\
%         (a) && (b)
%     \end{tabular}
%     \caption{Simulated environment featuring a UR5 manipulator with a flashlight end-effector in a room containing multiple rectangular windows.}% The targets are randomly generated outside the room, requiring the robot to find a window through which it can illuminate the target.}
%     \label{fig:simulated:experiments}
% \end{figure}

\begin{figure}
\vspace{0.3cm}
\centering
\begin{tabular}{cc}
    \includegraphics[height=3.2cm]{figs/new_envs/door_task.png} &  \includegraphics[height=3.2cm]{figs/new_envs/windows_task.png}\\
    % \includegraphics[trim={50mm 0mm 60mm 60mm}, clip, height=3.5cm]{figs/new_envs/double_rooms_2_1.png} &  \includegraphics[trim={20mm 0mm 20mm 0mm}, clip, height=3.5cm]{figs/new_envs/double_rooms_2_4.png}\\
\end{tabular}
\caption{\new{Simulated (left) door and (right) windows environments with a 9-DOF robot moving to illuminate a sphere.}}% described in Section \ref{sec:experiments:sim}. The experiment requires a flashlight-wielding mobile UR5 to illuminate targets in the open containers in the right room, and on the shelves in the door-less right room.}}
\label{fig:simulated:experiments}
\vspace{-0.7cm}
\end{figure}
%%%%%%%%%%%%%%%%%%%%%%%%%%%%%

Each of the five tested algorithms was invoked for \new{200 trials with a time limit of 120 seconds each}. The \rrt variants, \rrt and \visrrt, were given a fixed start configuration and no information was saved between runs of the experiment. \prm variants, \visprm, \vir, and \prm itself, were given randomized valid start configurations, \new{excluding configurations inside door-less room in order to assure task feasibility}, and maintained the aggregated roadmap throughout the entire experiment. Runtime statistics and success rates were recorded, and are presented in Table \ref{tab:sim:exp}. The results show a clear success rate and runtime advantage to \visrrt and \visprm over the compared algorithms. \new{The separation of the two tasks demonstrates that the narrow passage in the door room favors the \prm variants, excluding the benchmark \prm. We attribute \prm's poor performance to the small probability of sampling a manipulator configuration illuminating inside the containers.}

Figure \ref{fig:results:amortized} shows the improvement in runtime of \prm variants as the number of queries increases. These trends suggest that the aggregated statistics in Table \ref{tab:sim:exp} obscure the superior performance of \visprm, which becomes significantly more pronounced as the roadmap grows. \new{Note that the runtime average includes the pre-processing time for building the \vi-tree and \vir's workspace clustering.} \remove{The runtime distributions for \rrt and \prm variants are illustrated in Figure \ref{fig:results:success:rate}.}For \prm, the minimal runtime deviation indicates that success is mostly limited to queries immediately answered by the current roadmap. %The lack of data points at higher time intervals suggests that if a solution is not found quickly, the baseline rarely succeeds through further sampling.



% ---------------------------------------------------

\subsection{Real-world demonstrations}

The proposed \tvmp algorithms were further validated on a physical UR5 manipulator. The robotic arm was situated in a constrained workspace featuring a table with a three-tiered shelving unit constructed from nested boxes. A coordinate grid was overlaid on the table and each shelf level to facilitate consistent positioning. The experimental task required the illumination of two target objects, a mug and a pen holder, placed at random locations. Examples of robot and objects positioning can be seen in Figures \ref{fig:tvmp}b and \ref{fig:physical:experiments:two:2}. A ZED 2 depth camera mounted on the ceiling tracked the 3D coordinates of these objects in real-time. Before each trial, the algorithm computed a bounding sphere encompassing both targets to define the visibility target for the \tvmp solver. 

%%%%%%%%%%%%%%%%%%%%%%%%%%%%%
\begin{table}
\vspace{0.2cm}
    \centering
    \caption{Performance results for the simulated experiments of the \new{window and door rooms}.} 
    % \new{Success rates are shown for all experiments and also separated by task (windows or door room).}}
    
    % Adjust the values in the curly braces below to change individual column widths
    \begin{tabular}{L{0.9cm} C{0.3cm} C{0.5cm} C{0.3cm} C{1.7cm} C{0.4cm} C{0.4cm} C{0.6cm}}
        \toprule

        \multirow{2}{*}{Algorithm} & \multicolumn{3}{c}{Success rate (\%)} & \multicolumn{4}{c}{Comp. time (s)} \\
        \cmidrule(lr){2-4} \cmidrule(lr){5-8} 
        & All & \new{Win.} & \new{Door} & Mean & Med. & Min. & Max. \\
        \midrule
        \visrrt & 64.0 & 90.4 & 29.1 & 6.61 $\pm$ 12.56  & 1.48 & 0.09 & 68.21 \\
        \rrt    & 37.5 & 62.3 & 4.7  & 14.94 $\pm$ 21.61 & 5.82 & 0.09 & 119.26 \\
        
        \midrule
        
        \visprm & 87.0 & 87.7 & 86.0 & 9.12 $\pm$ 15.15 & 2.48 & 0.11 & 94.81 \\
        \vir     & 69.0 & 75.4 & 60.5 & 9.80 $\pm$ 15.84  & 4.30 & 0.25 & 97.23 \\
        \prm    & 44.0 & 77.2 & 0.0  & 4.97 $\pm$ 9.53  & 1.61 & 0.28 & 71.78 \\
        \bottomrule
    \end{tabular}
    \label{tab:sim:exp}
    \vspace{-0.6cm}
\end{table}
%%%%%%%%%%%%%%%%%%%%%%%%%%%%%
%%%%%%%%%%%%%%%%%%%%%%%%%%%%%
\begin{figure}[h]
    \centering
    \includegraphics[width=\linewidth]{figs/new_results/overall_amortized_runtime.png}
    \vspace{-0.7cm}
    \caption{\new{Amortized runtime across 200 successful and failed trials, including \visprm, \vir  and \prm overhead. Each point denotes the cumulative runtime over consecutive trials.}
    % \new{Amortized runtime of all runs, i.e., both successful and failed, and including pre-processing time for \visprm (blue), \vir (orange), and \prm (green).  Specifically, the $y$-value at any given $x$-coordinate $X$, represents the average runtime of the first $X$ runs out of 200}.
    }
    \label{fig:results:amortized}
    \vspace{-0.3cm}
\end{figure}
%%%%%%%%%%%%%%%%%%%%%%%%%%%%%

While the environment constrained the robot's motions, the ZED 2 setup prohibited creating more occlusion. To compensate and increase task difficulty, we restricted the flashlight’s \fov to a narrow beam of $\gamma=10^\circ$ and enforced a 60-second time limit per query. We compared all four algorithms under these settings, while also including a $30^\circ$ \fov variant of \prm as an additional baseline. Each algorithm was evaluated over 20 trials. The \rrt variants were initialized from a fixed starting configuration. In contrast, the \prm variants planned from the robot’s current state toward the target sets, with roadmaps incrementally aggregated throughout the duration of the experiment.

\remove{
%%%%%%%%%%%%%%%%%%%%%%%%%%%%%
\begin{figure}[h]
    \centering
    \includegraphics[width=\linewidth]{figs/new_results/overall_scatter_and_success_rate.pdf}
    % \vspace{-0.3cm}
    \caption{Comparative runtime distributions (seconds) for successful trials in simulation: (left) \rrt vs. \visrrt, and (right) \visprm vs. \new{\vir and} \prm . Individual data points denote the runtime of single successful executions, with red circles indicating mean values and orange horizontal lines representing medians.}
    \label{fig:results:success:rate}
\end{figure}
%%%%%%%%%%%%%%%%%%%%%%%%%%%%%
}
% %%%%%%%%%%%%%%%%%%%%%%%%%%%%%
% \begin{figure}[h]
%     \centering
%     \includegraphics[height=3.8cm]{figs/physical_experiment_2_os.png}
%     % \vspace{-0.3cm}
%     \caption{Experimental setup featuring the physical manipulator in a constrained workspace, tasked with illuminating two red target objects.}
%     \label{fig:physical:experiments:two:1}
% \end{figure}
% %%%%%%%%%%%%%%%%%%%%%%%%%%%%%
% %%%%%%%%%%%%%%%%%%%%%%%%%%%%%
% \begin{figure}[h]
%     \centering
%     \includegraphics[height=3.7cm]{figs/illuminate_ball_o.png}
%     % \vspace{-0.3cm}
%     \caption{A secondary test scenario requiring the robot to achieve illumination of a ball target (marked in red) occluded by a wall.}
%     \label{fig:physical:experiments:two:2}
%     \vspace{-0.3cm}
% \end{figure}
% %%%%%%%%%%%%%%%%%%%%%%%%%%%%%
% %%%%%%%%%%%%%%%%%%%%%%%%%%%%%
\begin{figure}[h]
    \vspace{0.3cm}
    \centering
    \begin{tabular}{cc}
        \includegraphics[height=3cm]{figs/physical_experiment_2_os.png} & \includegraphics[height=3cm]{figs/illuminate_ball_o_2.png} \\
        (a) & (b)
    \end{tabular}
    % \vspace{-0.3cm}
    \caption{Experiments of (a) a manipulator illuminating two red targets within a constrained workspace, and (b) a test scenario requiring illumination of a wall-occluded ball.}
    \label{fig:physical:experiments:two:2}
\end{figure}
% %%%%%%%%%%%%%%%%%%%%%%%%%%%%%

%%%%%%%%%%%%%%%%%%%%%%%%%%%%%
% \begin{figure}[t!]
%     \centering
%     \begin{tabular}{ccc}
%          \includegraphics[height=3.7cm]{figs/physical_experiment_2.png} &~~~ & \includegraphics[height=3.7cm]{figs/illuminate_ball.png}\\
%          (a) && (b)
%     \end{tabular}       
%     \caption{(a) Experimental setup featuring the physical manipulator in a constrained workspace, tasked with illuminating two red target objects. (b) A secondary test scenario requiring the robot to achieve illumination of a ball target (marked in red) occluded by a wall.}
%     \label{fig:physical:experiments:two}
% \end{figure}
% \begin{figure*}%[!t]
% \centering
% % \subfloat[]{\includegraphics[height=3.8cm]{figs/physical_experiment_2.jpg}%
% % \label{fig_phys_exp_a}}
% % \hfil
% % \subfloat[]{\includegraphics[height=3.8cm]{figs/illuminate_ball.png}%
% % \label{fig_phys_exp_b}}
% \begin{tabular}{ccc}
%          \includegraphics[height=3.7cm]{figs/physical_experiment_2.png} &~~~ & \includegraphics[height=3.7cm]{figs/illuminate_ball.png}\\
%          % (a) && (b)
%     \end{tabular} 
% \caption{(a) Experimental setup featuring the physical manipulator in a constrained workspace, tasked with illuminating two red target objects. (b) A secondary test scenario requiring the robot to achieve illumination of a ball target (marked in red) occluded by a wall.}
% \label{fig:physical:experiments:two}
% \vspace{-0.3cm}
% \end{figure*}
%%%%%%%%%%%%%%%%%%%%%%%%%%%%%

%%%%%%%%%%%%%%%%%%%%%%%%%%%%%
% \begin{figure}
%     \centering
%     \includegraphics[width=0.5\linewidth]{figs/physical_experiment_2.png}
%     \caption{The setup of the real-world experiment. The robot is placed in a constrained environment and tasked with illuminating two red objects.}
%     \label{fig:physical:experiment}
% \end{figure}
%%%%%%%%%%%%%%%%%%%%%%%%%%%%%

\new{Query success required generating a valid path and physically executing it to illuminate both target objects within a strict 60-second time limit, with results summarized in Table \ref{tab:physical:exp}. Despite the highly constrained 60-second budget and a narrow $10^\circ$ \fov{} beam, \visprm and \visrrt significantly outperformed the baseline benchmarks, resolving approximately half of all queries. The baseline \prm devoted more computational overhead to nearest-neighbor searches and edge validation, adding fewer nodes and examining fewer \fov{}s than \rrt. Furthermore, the spatial concentration of targets on the table surface favored \visprm and \visrrt, as configurations sampled by the \vi-tree to observe one target were highly reusable for others. Finally, to match the coverage efficiency inherently provided by our visibility-aware sampling, the baseline \prm required increasing the FOV angle to $\gamma=30^\circ$—a threefold increase in angle yielding a roughly 9-times larger \fov area. Another demonstration of illuminating an occluded ball target is shown in Figure \ref{fig:physical:experiments:two:2}b and in the supplementary video.}

\vspace{-0.1cm}

% Query success was defined by the generation of a valid path and its subsequent physical execution to successfully illuminate both target objects, with results summarized in Table \ref{tab:physical:exp}. \visprm and \visrrt significantly outperformed the baseline benchmarks, resolving approximately half of all queries. The performance disparity between the baseline \rrt and \prm may be attributed to the limited trial count, yet it remains consistent with the total number of nodes added to their respective graphs. Specifically, the baseline \prm devoted more computational overhead to nearest-neighbor searches and edge validation, thereby reducing the total number of \fov{}s examined compared to \rrt. Furthermore, the spatial concentration of targets on the table surface, defined by a fixed $z$-value and narrow $x$ and $y$ intervals, likely favored \visprm over \visrrt, as configurations sampled by the \vi-tree to observe one target were highly likely to be reused for subsequent targets. Finally, the \prm variant with $\gamma=30^\circ$ achieved success rates comparable to our proposed methods, providing critical intuition: the baseline required an increase in angle by a factor of 3, yielding a roughly 9-times increase in \fov area, to match the coverage efficiency inherently provided by our visibility-aware sampling. Another demonstration of illuminating a ball target occluded by a wall is seen in Figure \ref{fig:physical:experiments:two:2}. 

%%%%%%%%%%%%%%%%%%%%%%%%%%%%%
% \begin{table}[t!]
%     \centering
%     \caption{Success rates for real robot experiments}
%     \begin{tabular}{lcc}
%         \toprule
%         Algorithm &~~~~& Success Rate (\%) \\
%         \midrule
%         \prm ($\gamma=10^\circ$)    && 0  \\
%         \prm ($\gamma=30^\circ$)    && 45 \\
%         \visprm                     && 50 \\
%         \rrt                        && 10 \\
%         \visrrt                     && 40 \\
%         \bottomrule
%     \end{tabular}
%     \label{tab:physical:exp}
% \end{table}
%%%%%%%%%%%%%%%%%%%%%%%%%%%%%

%%%%%%%%%%%%%%%%%%%%%%%%%%%%%%%%
\begin{table}
    \centering
    \caption{Success rates for real robot experiments}
    \footnotesize
    \setlength{\tabcolsep}{8pt} % <--- Increases space between columns
    \begin{tabular}{lccccc}
        \toprule
        Algorithm & \prm & \prm & \visprm & \rrt & \visrrt \\
        & ($10^\circ$) & ($30^\circ$) & & & \\
        \midrule
        Suc. Rate (\%) & 0 & 45 & 50 & 10 & 40 \\
        \bottomrule
    \end{tabular}
    \label{tab:physical:exp}
    \vspace{-0.4cm}
\end{table}
%%%%%%%%%%%%%%%%%%%%%%%%%%%%%%%%%


% ------------------ v OLD v -------------------------
% ---------------------------------------------------
% \subsection{Experimental setup}
% \label{sec:experiments:setup}




% \paragraph{Robot:} All of our experiments, simulated and physical, use a UR5 manipulator fitted with an adjustable focus flashlight connected to its end-effector. The manipulator is mounted a Clearpath Robotics Husky A200, a 4-wheel mobile platform, but \textbf{the mobile platform was not moved in any of the experiments}. See Figure \ref{fig:physical:experiment}. We assign a range of 3 meters to the flashlight, and a beam half-angle of $\frac{\pi}{12}$ ($\approx 15$ degrees) in the simulated experiment, and $\frac{\pi}{36}$ ($\approx 5$ degrees) in the physical experiment, except for one variant of \prm (See Section \ref{sec:experiments:results:sim}). 

% \paragraph{Task randomization:}
% Random variables determining tasks, that is, start configurations and targets, were instantiated before each of the two experiments and used for all algorithms. In-algorithm randomness, that is, all sampling behaviors, was \textbf{not} fixed in the same way.   

% \paragraph{Task admissibility:}
% Randomly generated targets were not guaranteed to give rise to admissible \tvmp problems. that is, for a randomly generated instance $(r,E,s,P_S)$, a configuration of $r$ that illuminates the desired fraction of the target set might not exist. To mitigate this, \textbf{All reported runtimes are of successful runs only}.

% \begin{figure}
%     \centering
%     \includegraphics[width=\linewidth]{figs/husky.jpg}
%     \caption{The robot used for the experiments - a Clearpath Robotics Husky A-200 mounted with a UR5 with a flashlight end-effector}
%     \label{fig:husky}
% \end{figure}

% \subsubsection{simulated experiment setup}

% In the simulated experiments, shown in Figure \ref{fig:simulated:experiments}, the robot is placed at the center of a 4-wall room, with windows of varying sizes and shapes in different parts of the wall. The robot is tasked with illuminating a sphere with center chosen outside of the walls uniformly at random within a range of $2.5$ meters from the robot (accommodating the $3$ meter flashlight range), and with radius between $3$ and $30$ centimeters. 


% \begin{figure}[t!]
%     \centering
%     \begin{subfigure}{0.47\linewidth}
%     \centering
%         \includegraphics[width=\linewidth]{walls_1.png}
%         \caption{}
%         \label{fig:simulated:experiments:walls}
%     \end{subfigure}
%     \hfill
%     \begin{subfigure}{0.47\linewidth}
%     \centering
%         \includegraphics[width=\linewidth]{walls_3.png}
%         \caption{}
%         \label{fig:simulated:experiments:walls:2}
%     \end{subfigure}
%     \caption{Images from the simulated experiments described in Section \ref{sec:experiments}. The mounted UR5 is equipped with a flashlight and positioned in a room with various rectangular windows. The targets are randomly generated outside the room, requiring the robot to find a window through which it can illuminate the target.}
%     \label{fig:simulated:experiments}
% \end{figure}

% Each of the four algorithms \prm, \visprm, \rrt, and \visrrt were invoked 100 times. The \rrt variants were given a fixed start configuration and no information was saved between runs of the experiment. \prm variants were given randomized valid start configurations, and maintained the aggregated roadmap throughout the entire experiment. Runtime statistics and success rates were recorded, and are presented in Section \ref{sec:experiments:results:sim}.


% \subsubsection{physical experiment setup}

% In the physical experiment, shown in Figure \ref{fig:physical:experiment}, the robot was placed in a constrained space, next to a table with nested boxes forming a three level shelf. The table and every level of the shelf were overlayed with a grid. The \tvmp task required the robot to illuminate two objects, a red mug and a red pen holder, placed at random locations. Tasks were generated by uniformly sampling the location of the first object, and uniformly sampling a location for the second object within distance 2, including diagonals, from the first.

% A Zed2 depth camera was placed above the workspace, and recorded the 3D location of the objects. Upon receiving the planning command an enclosing sphere of the two objects was computed and the\tvmp algorithm was called.

% Each of the four algorithms \prm, \visprm, \rrt, and \visrrt were invoked 20 times. The \rrt variants started from a fixed start configuration and no information was saved between runs of the experiment. \prm variants were only given the target sets, and planned motions starting from the robot's current configuration. Roadmaps were aggregated throughout the entire experiment. 

% While the environment constrained the robot's motions, the Zed setup prohibited us from creating more occlusion. We chose to increase the queries' difficulty by using a relatively narrow \fov, i.e., light beam, and a time cap of 60 seconds. All four algorithms, \prm, \visprm, \rrt, and \visrrt were compared with the \fov half-angle set to $\pi/36$, alongside an additional variant of \prm with a half-angle of $\pi/12$. Success rates were recorded, and are presented in Section \ref{sec:experiments:results:physical}.

% \begin{figure}
%     \centering
%     \includegraphics[width=0.7\linewidth]{figs/physical_experiment_2.jpg}
%     \caption{The setup of the physical experiment. The robot is placed in a constrained environment and tasked with illuminating a target set. The target set is observed by an overhead Zed2 depth camera which sends its location to the robot.}
%     \label{fig:physical:experiment}
% \end{figure}


% \subsection{Implementation details}
% Here we give a review of our implementation, detailing design and parameter choices used for the experiments. All of our code is available at \url{https://github.com/StavAshur/visual_based_planning}


% \subsubsection{System details} All experiments were conducted on a desktop PC with Intel i7-14700F CPU, 64 GB of RAM, NVIDIA GeForce RTX 4090 GPU. To maintain compatibility with the Husky A-200 robot, all software components used were chosen for compatibility with ROS Melodic running within an Ubuntu 18.04 and built on an Ubuntu 18.04 Distrobox container hosted on an Ubuntu 24.04 operating system.





% \subsection{Results and discussion}
% \label{sec:experiments:results}

% \subsubsection{Simulated experiment}
% \label{sec:experiments:results:sim}
% The simulated experiment's results, taking only successful runs into account,  are summed up in Table \ref{tab:sim:exp}. These results show a clear success rate advantage to \visrrt and \visprm over the benchmark algorithms solving $\sim30\% $ more queries. \visrrt also shows much faster runtimes, with $\sim40\% $ faster mean runtime, and $\sim\times4$ speed up in the median case.

% The max. value results paired with the discrepancy in success rates, might indicate that the benchmarks' were unable to solve the hardest cases. While the benchmarks were able to solve some instances that were not solved by \visrrt and \visprm, the time required to solve these was relatively short. The distribution of runtimes can be seen in Figure \ref{fig:results:success:rate}.

% In particular, the statistics of \prm show a very low runtime standard deviation, and the scatter plot of Figure \ref{fig:results:success:rate} corroborates the intuition that very few queries that are not answered immediately, presumably by scanning configurations already in the roadmap, are successfully answered by \prm.

% Figure \ref{fig:results:amortized} rovement in success rate and runtime of \prm and \visprm as the number of past queries grows. This shows that similarity between the runtime and success rate statistics of \visrrt and \visprm in Table \ref{tab:sim:exp} hides better \visprm performance after the roadmap has grown sufficiently.  



% \begin{table}[t!]
%     \centering
%     \caption{Simulated experiment success rate and runtime statistics}
%     \begin{tabular}{ l c ccccc}
%         \toprule
%         Alg. & Succ. Rate & Mean (s) & Median (s) & Min (s) & Max (s) & Std Dev (s) \\
%         \midrule
%         \prm & 56/100 & 3.74 & 1.33 & 0.524 & 24.04 & 5.22 \\
%         \visprm & 77/100 & 8.50 & 1.76 & 0.265 & 89.40 & 16.74 \\
%         \rrt & 51/100 & 13.53 & 8.14 & 0.059 & 44.94 & 13.14 \\
%         \visrrt & 78/100 & 8.29 & 2.33 & 0.087 & 86.92 & 14.93 \\
%         \bottomrule
%     \end{tabular}
%     \label{tab:sim:exp}
% \end{table}


% \begin{figure}
%     \centering
%     \includegraphics[width=\linewidth]{figs/cumulative_succ_runtime_prm.pdf}
%     \caption{Cumulative success rate (left) and runtime (right) results for the 100 \prm (orange) and \visprm (blue) simulation experiments. The $y$-value represents the success rate/runtime average in all previous queries. Note that the $x$-axis, marking the number of queries included in the calculation, starts at $10$.}
%     \label{fig:results:cumulative}
% \end{figure}

% \begin{figure}
%     \centering
%     \includegraphics[width=\linewidth]{figs/succ_runtime_rrt_prm.pdf}
%     \caption{Runtime results (in seconds) for \textbf{successful runs} in the simulated experiments. Left: \rrt vs. \visrrt. Right: \prm vs. \visprm. Every point represents the runtime of a single successful run, the bar shows the mean runtime, and the red line marks the median runtime.}
%     \label{fig:results:success:rate}
% \end{figure}




% \subsubsection{Physical experiment}
% \label{sec:experiments:results:physical}

% The physical experiment's results are summed up in Table \ref{tab:physical:exp}. the two \prm variants, \prm10 and \prm30 are \prm{}s with an \fov with a 10 degrees and 30 degrees angles respectively. Once more, \visprm and \visrrt achieve a much higher success rate than the benchmarks, solving roughly half of the queries. 

% While the different success rates of \rrt and \prm might be an artifact of the low number of experiments, they are also consistent with the amount of overall configurations added to the graph. \prm spends considerably more time searching for nearest neighbors and validating edges, reducing the overall number of \fov{}s examined by \prm compared to \rrt. 

% Similarly, we hypothesize that the concentration of targets in a single $z$-value and small $x$ and $y$ intervals, namely, the surface of the table, favored \visprm over \visrrt. Configurations sampled using the \vi tree to observe a target were likely to be reused as targets due to this property of the target distribution.

% The success rate of \prm30, similar to those of \visprm and \visrrt, provides more intuition regarding the advantage of our methods over the benchmark. The $\times3$ larger angle, resulting in a $\approx\times9$ larger \fov was required for the coverage of \prm to allow this success rate.

% \begin{table}[t!]
%     \centering
%     \caption{physical experiment success rate}
%     \begin{tabular}{|l|c|}
%         \hline
%         Alg. & Succ. Rate \\
%         \hline
%         \prm10 & 0/20 \\
%         \prm30 & 9/20 \\
%         \visprm & 10/20 \\
%         \rrt & 2/20 \\
%         \visrrt & 8/20 \\
%         \hline
%     \end{tabular}
%     \label{tab:physical:exp}
% \end{table}

